\section{2016 Analysis \& Differential Equations}

\subsection{Individual}

\begin{exercise}
\label{analysis:2016:1}
Suppose that $F$ is continuous on $[a,b]$, $F'(x)$ exists for every $x\in(a,b)$, and $F'(x)$ is integrable. Prove that $F$ is absolutely continuous and
\[
F(b) - F(a) = \int_a^b F'(x)\,dx.
\]
\end{exercise}

\begin{solution}
\textbf{Proof.}
We use the standard contest interpretation that $F'$ is Riemann integrable on $[a,b]$. Then $F'$ is bounded; choose $M$ such that $|F'(x)|\leq M$ on $(a,b)$. By the mean value theorem,
\[
|F(y)-F(x)|\leq M|y-x|
\]
for all $x,y\in[a,b]$. Hence $F$ is Lipschitz, and in particular absolutely continuous.

It remains to prove the integral identity. Let
\[
P:\quad a=x_0<x_1<\cdots <x_N=b
\]
be any partition. For each interval $[x_{j-1},x_j]$, the mean value theorem gives a point $\xi_j\in(x_{j-1},x_j)$ such that
\[
F(x_j)-F(x_{j-1})=F'(\xi_j)(x_j-x_{j-1}).
\]
Summing over $j$,
\[
F(b)-F(a)=\sum_{j=1}^N F'(\xi_j)(x_j-x_{j-1}).
\]
The right-hand side is a Riemann sum for $F'$. Since $F'$ is Riemann integrable, taking partitions with mesh tending to zero gives
\[
F(b)-F(a)=\int_a^b F'(x)\,dx.
\]
This also agrees with the absolute-continuity representation
\[
F(x)=F(a)+\int_a^x F'(t)\,dt.
\]
\end{solution}

\begin{exercise}
\label{analysis:2016:2}
Suppose that $f$ is integrable on $\mathbb{R}^n$, let $K_{\delta}(x) = \delta^{-\frac{n}{2}}e^{-\frac{\pi|x|^2}{\delta}}$ for each $\delta>0$. Prove that the convolution
\[
(f*K_{\delta})(x) = \int_{\mathbb{R}^n} f(x-y)K_{\delta}(y)\,dy
\]
is integrable and $\| (f*K_{\delta}) - f\|_{L^1(\mathbb{R}^n)} \to 0$ as $\delta\to 0$.
\end{exercise}

\begin{solution}
\textbf{Prerequisites:}
\begin{itemize}
    \item \textbf{Convolution:} The operation $(f * g)(x) = \int_{\mathbb{R}^n} f(x-y)g(y) dy$.
    \item \textbf{Young's Inequality (for $L^1$):} $\|f * g\|_1 \leq \|f\|_1 \|g\|_1$.
    \item \textbf{Gaussian Integral:} $\int_{-\infty}^\infty e^{-ax^2} dx = \sqrt{\pi/a}$. For $\mathbb{R}^n$, $\int_{\mathbb{R}^n} e^{-\pi |x|^2} dx = 1$.
    \item \textbf{Continuity of Translation in $L^1$:} For any $f \in L^1(\mathbb{R}^n)$, the mapping $y \mapsto f(\cdot - y)$ is continuous from $\mathbb{R}^n$ to $L^1(\mathbb{R}^n)$. That is, $\lim_{y \to 0} \|f(\cdot - y) - f(\cdot)\|_1 = 0$.
\end{itemize}

\textbf{Intuition:}
The function $K_\delta(x)$ is a "mollifier" or an "approximation to the identity." As $\delta$ becomes very small, $K_\delta$ becomes a very sharp peak at the origin, while the area under the curve remains exactly 1. Convolving $f$ with $K_\delta$ is like taking a weighted average of $f$ around each point $x$. Since the weight is concentrated near the origin when $\delta$ is small, the result $(f * K_\delta)(x)$ should be very close to $f(x)$.

\textbf{Steps:}

\textbf{Part 1: Proof of Integrability}
\begin{enumerate}
    \item \textbf{Apply Young's Inequality:}
    By Young's inequality for convolutions, the $L^1$ norm of the product is bounded by the product of the $L^1$ norms:
    $\|f * K_\delta\|_1 \leq \|f\|_1 \|K_\delta\|_1$.

    \item \textbf{Compute the norm of the Kernel:}
    We need to calculate $\|K_\delta\|_1 = \int_{\mathbb{R}^n} \delta^{-n/2} e^{-\pi |x|^2 / \delta} dx$.
    Use the substitution $u = x / \sqrt{\delta}$, which implies $dx = \delta^{n/2} du$.
    The integral becomes:
    $\int_{\mathbb{R}^n} \delta^{-n/2} e^{-\pi |u|^2} (\delta^{n/2} du) = \int_{\mathbb{R}^n} e^{-\pi |u|^2} du = 1$.
    Since $\|f\|_1$ is finite and $\|K_\delta\|_1 = 1$, the convolution $\|f * K_\delta\|_1$ is finite, meaning it is integrable.
\end{enumerate}

\textbf{Part 2: Proof of $L^1$ Convergence}
\begin{enumerate}
    \item \textbf{Setup the difference integral:}
    Since $\int K_\delta(y) dy = 1$, we can write $f(x) = \int f(x) K_\delta(y) dy$.
    Then, $\|f * K_\delta - f\|_1 = \int_{\mathbb{R}^n} | \int_{\mathbb{R}^n} (f(x-y) - f(x)) K_\delta(y) dy | dx$.
    By Fubini's Theorem (swapping the order of integration):
    $\|f * K_\delta - f\|_1 \leq \int_{\mathbb{R}^n} K_\delta(y) \left( \int_{\mathbb{R}^n} |f(x-y) - f(x)| dx \right) dy$.

    \item \textbf{Analyze the inner integral:}
    Let $\omega_f(y) = \int_{\mathbb{R}^n} |f(x-y) - f(x)| dx$. This is the $L^1$ "modulus of continuity."
    We know two things:
    (a) $\lim_{y \to 0} \omega_f(y) = 0$.
    (b) $\omega_f(y) \leq \|f(\cdot-y)\|_1 + \|f\|_1 = 2\|f\|_1$ (boundedness).

    \item \textbf{Split the domain of $y$:}
    For any $\epsilon > 0$, choose $\eta > 0$ such that $|y| < \eta \implies \omega_f(y) < \epsilon$.
    Split the outer integral into two parts: $|y| < \eta$ and $|y| \geq \eta$.
    $I_1 = \int_{|y|<\eta} K_\delta(y) \omega_f(y) dy < \epsilon \int K_\delta(y) dy = \epsilon$.
    $I_2 = \int_{|y|\geq \eta} K_\delta(y) \omega_f(y) dy \leq 2\|f\|_1 \int_{|y|\geq \eta} K_\delta(y) dy$.

    \item \textbf{Take the limit as $\delta \to 0$:}
    As $\delta \to 0$, the mass of $K_\delta$ concentrates at the origin.
    Specifically, $\int_{|y|\geq \eta} K_\delta(y) dy = \int_{|u| \geq \eta/\sqrt{\delta}} e^{-\pi |u|^2} du$.
    As $\delta \to 0$, the lower limit $\eta/\sqrt{\delta} \to \infty$, so the integral goes to 0.
    Thus, $\limsup_{\delta \to 0} \|f * K_\delta - f\|_1 \leq \epsilon$. Since $\epsilon$ is arbitrary, the limit is 0.
\end{enumerate}
\end{solution}

\begin{exercise}
\label{analysis:2016:3}
Prove that a bounded function on interval $I=[a,b]$ is Riemann integrable if and only if its set of discontinuities has measure zero.
\end{exercise}

\begin{solution}
\textbf{Prerequisites:}
\begin{itemize}
    \item \textbf{Riemann Integrability:} A bounded function $f$ on $[a, b]$ is Riemann integrable if for every $\epsilon > 0$, there exists a partition $P$ such that the difference between the upper and lower Darboux sums $U(P, f) - L(P, f) < \epsilon$.
    \item \textbf{Oscillation:} The oscillation of $f$ at a point $x$ is $\omega(f, x) = \lim_{\delta \to 0} \sup \{|f(y) - f(z)| : y, z \in (x-\delta, x+\delta) \cap [a, b]\}$. A function is continuous at $x$ if and only if $\omega(f, x) = 0$.
    \item \textbf{Lebesgue Measure Zero:} A set $S \subset \mathbb{R}$ has measure zero if for every $\epsilon > 0$, it can be covered by a countable collection of open intervals whose total length is less than $\epsilon$.
\end{itemize}

\textbf{Intuition:}
Riemann integration works by approximating a function with rectangles. If a function has too many discontinuities, the heights of these rectangles (the "gap" between the maximum and minimum values in a subinterval) won't shrink to zero even as the widths do. Lebesgue's criterion tells us that as long as the "bad" points (discontinuities) are few enough to be ignored in terms of length (measure zero), the function remains integrable.

\textbf{Steps:}

\textbf{Step 1: Characterize the set of discontinuities}
Let $D$ be the set of points where $f$ is discontinuous. As noted, $x \in D \iff \omega(f, x) > 0$.
We can write $D$ as a countable union: $D = \bigcup_{k=1}^\infty D_{1/k}$, where $D_\epsilon = \{x \in [a, b] : \omega(f, x) \geq \epsilon\}$.
It is a known property that $D_\epsilon$ is always a \textit{closed} and \textit{bounded} (hence compact) set.

\textbf{Step 2: Proof that $m(D)=0 \implies$ Riemann Integrability}
\begin{enumerate}
    \item Assume $m(D) = 0$. This implies $m(D_\epsilon) = 0$ for all $\epsilon > 0$.
    \item Fix $\epsilon > 0$. Since $m(D_\epsilon) = 0$, we can cover $D_\epsilon$ by a finite number of open intervals $(a_i, b_i)$ with total length $\sum (b_i - a_i) < \epsilon$. (Finite because $D_\epsilon$ is compact).
    \item The remaining part of $[a, b] \setminus \bigcup (a_i, b_i)$ is a collection of closed intervals. On these intervals, the oscillation of $f$ is strictly less than $\epsilon$.
    \item We can create a partition $P$ using the endpoints of these intervals.
    \item In the sum $U(P, f) - L(P, f) = \sum (M_j - m_j) \Delta x_j$:
       \begin{itemize}
           \item For subintervals covering $D_\epsilon$, the total contribution is $\leq (M - m) \cdot \sum \Delta x_j < (M - m)\epsilon$, where $M, m$ are the global bounds of $f$.
           \item For the other subintervals, the contribution is $\leq \epsilon \cdot (b - a)$.
       \end{itemize}
    \item Thus, $U - L \leq \epsilon(M - m + b - a)$. Since $\epsilon$ is arbitrary, $f$ is Riemann integrable.
\end{enumerate}

\textbf{Step 3: Proof that Riemann Integrability $\implies m(D)=0$}
\begin{enumerate}
    \item Assume $f$ is Riemann integrable. Then for any $\eta > 0$, there is a partition $P$ such that $U(P, f) - L(P, f) < \eta$.
    \item For a fixed $\epsilon > 0$, consider the set $D_\epsilon$. If a subinterval $I_j = [x_{j-1}, x_j]$ of the partition contains a point of $D_\epsilon$ in its interior, then $M_j - m_j \geq \epsilon$.
    \item Let $J$ be the set of indices of subintervals that contain points of $D_\epsilon$.
    \item Then $\eta > U - L \geq \sum_{j \in J} (M_j - m_j) \Delta x_j \geq \epsilon \sum_{j \in J} \Delta x_j$.
    \item This means the total length of subintervals covering $D_\epsilon$ (excluding partition points, which are finite) is $\sum_{j \in J} \Delta x_j < \eta / \epsilon$.
    \item Since $\eta$ can be chosen arbitrarily small, $m(D_\epsilon) = 0$.
    \item Since $D = \bigcup_{k=1}^\infty D_{1/k}$ and each $D_{1/k}$ has measure zero, their union $D$ also has measure zero.
\end{enumerate}
\end{solution}

\begin{exercise}
\label{analysis:2016:4}
\begin{itemize}
    \item[1)] Let $w=\frac{1}{2}(z+\frac{1}{z})$. Show that it maps $\{|z|>1\}$ to $\mathbb{C}\setminus[-1,1]$.
    \item[2)] Compute $\int_0^\infty \frac{\log x}{x^2-1} dx$.
\end{itemize}
\end{exercise}

\begin{solution}
\textbf{Prerequisites:}
\begin{itemize}
    \item \textbf{Complex Mapping:} Representing $z = x + iy$ as $re^{i\theta}$.
    \item \textbf{Ellipses:} The standard equation $u^2/a^2 + v^2/b^2 = 1$ with semi-axes $a$ and $b$.
    \item \textbf{Improper Integrals:} Handling $\log x$ and singularities at $x=1$.
    \item \textbf{Taylor Series:} $1/(1-r) = \sum r^n$ for $|r| < 1$.
    \item \textbf{Riemann Zeta Function:} $\sum_{k=1}^\infty \frac{1}{k^2} = \frac{\pi^2}{6}$ (Basel Problem).
\end{itemize}

\textbf{Intuition:}
\textbf{Part 1:} The map $w = \frac{1}{2}(z + 1/z)$ is known as the Zhukovsky transform. It "squashes" the complex plane. Think of circles $|z|=r$. For $r=1$, it maps the circle to the segment $[-1, 1]$ twice. For $r > 1$, it maps circles to larger and larger ellipses that eventually cover the entire plane as $r \to \infty$.

\textbf{Part 2:} The integral $\int_0^\infty \frac{\log x}{x^2-1} dx$ looks tricky because of the $x^2-1$ in the denominator. However, the $\log x$ also vanishes at $x=1$, making the singularity removable. The symmetry between $[0, 1]$ and $[1, \infty]$ under the map $x \to 1/x$ allows us to simplify the problem to an interval $[0, 1]$ where we can use series.

\textbf{Steps:}

\textbf{Part 1: The Mapping}
\begin{enumerate}
    \item Let $z = re^{i\theta}$ where $r > 1$. Then $1/z = \frac{1}{r}e^{-i\theta}$.
    \item Substitute into $w = u + iv$:
    $u + iv = \frac{1}{2}(re^{i\theta} + \frac{1}{r}e^{-i\theta}) = \frac{1}{2}(r(\cos\theta + i\sin\theta) + \frac{1}{r}(\cos\theta - i\sin\theta))$.
    \item Separate the real and imaginary parts:
    $u = \frac{1}{2}(r + 1/r)\cos\theta$ and $v = \frac{1}{2}(r - 1/r)\sin\theta$.
    \item Define $a = \frac{1}{2}(r + 1/r)$ and $b = \frac{1}{2}(r - 1/r)$. Then $\frac{u^2}{a^2} + \frac{v^2}{b^2} = \cos^2\theta + \sin^2\theta = 1$.
    \item This is an ellipse. Since $r > 1$, $a > 1$ and $b > 0$. As $r \to 1^+$, $a \to 1$ and $b \to 0$, so the ellipse collapses to the interval $[-1, 1]$.
    \item As $r \to \infty$, $a, b \to \infty$. Thus, the image of $\{|z| > 1\}$ is the entire plane $\mathbb{C}$ minus the "collapsed" starting segment $[-1, 1]$.
\end{enumerate}
\textbf{Part 2: The Integral}
\begin{enumerate}
    \item Split the integral: $I = \int_0^1 \frac{\log x}{x^2-1} dx + \int_1^\infty \frac{\log x}{x^2-1} dx$.
    \item In the second part, let $x = 1/t$. Then $dx = -1/t^2 dt$.
       $\int_1^\infty \frac{\log x}{x^2-1} dx = \int_1^0 \frac{\log(1/t)}{(1/t)^2-1} (-1/t^2) dt = \int_0^1 \frac{-\log t}{\frac{1-t^2}{t^2}} \frac{1}{t^2} dt = \int_0^1 \frac{\log t}{t^2-1} dt$.
    \item Thus, $I = 2 \int_0^1 \frac{\log x}{x^2-1} dx$.
    \item Note that $\frac{1}{x^2-1} = \frac{-1}{1-x^2} = -\sum_{n=0}^\infty x^{2n}$.
       $I = -2 \int_0^1 \log x \left( \sum_{n=0}^\infty x^{2n} \right) dx = 2 \sum_{n=0}^\infty \int_0^1 (-x^{2n} \log x) dx$.
    \item Using integration by parts (or the formula $\int_0^1 x^m \log x dx = \frac{-1}{(m+1)^2}$):
       $\int_0^1 x^{2n} \log x dx = \left[ \frac{x^{2n+1}}{2n+1}\log x \right]_0^1 - \int_0^1 \frac{x^{2n+1}}{2n+1} \frac{1}{x} dx = 0 - \frac{1}{(2n+1)^2}$.
    \item $I = 2 \sum_{n=0}^\infty \frac{1}{(2n+1)^2}$.
    \item We know $\sum_{k=1}^\infty \frac{1}{k^2} = \sum_{k \text{ even}} \frac{1}{k^2} + \sum_{k \text{ odd}} \frac{1}{k^2} = \frac{1}{4}\sum_{j=1}^\infty \frac{1}{j^2} + \sum_{n=0}^\infty \frac{1}{(2n+1)^2}$.
       So $\sum_{n=0}^\infty \frac{1}{(2n+1)^2} = (1 - 1/4) \zeta(2) = \frac{3}{4} \cdot \frac{\pi^2}{6} = \frac{\pi^2}{8}$.
    \item Therefore, $I = 2 \cdot \frac{\pi^2}{8} = \frac{\pi^2}{4}$.
\end{enumerate}
\end{solution}

\begin{exercise}
\label{analysis:2016:5}
Let $f$ be a doubly periodic meromorphic function. Prove that the number of zeros and poles are equal.
\end{exercise}

\begin{solution}
Choose a fundamental parallelogram
\[
Q=z_0+[0,1]+i[0,1]
\]
whose boundary contains no zero or pole of $f$; this is possible after a small translation because zeros and poles are discrete modulo the lattice. Count zeros and poles in $Q$ with multiplicity.

By the argument principle,
\[
Z-P=\frac{1}{2\pi i}\int_{\partial Q}\frac{f'(z)}{f(z)}\,dz.
\]
The logarithmic derivative $f'/f$ has the same periods as $f$, namely $1$ and $i$. Therefore the integrals over opposite sides of $\partial Q$ cancel: the two horizontal sides are translates by $i$ and are traversed in opposite directions, and the same is true for the vertical sides, which are translates by $1$. Hence
\[
\int_{\partial Q}\frac{f'(z)}{f(z)}\,dz=0.
\]
Thus $Z-P=0$, so the number of zeros equals the number of poles in a fundamental parallelogram. Equivalently, a nonzero doubly periodic meromorphic function on the torus $\mathbb C/(\mathbb Z+i\mathbb Z)$ has divisor of degree zero.
\end{solution}

\begin{exercise}
\label{analysis:2016:6}
Prove that the spectrum of a bounded self-adjoint operator is a bounded closed subset of the real line.
\end{exercise}

\begin{solution}
\textbf{Boundedness and closedness.}
For every bounded operator $A$,
\[
\sigma(A)\subset \{\lambda\in\mathbb C:|\lambda|\leq \|A\|\}.
\]
Indeed, if $|\lambda|>\|A\|$, then
\[
A-\lambda I=-\lambda\left(I-\frac{A}{\lambda}\right)
\]
is invertible by the Neumann series. The resolvent set is open because invertibility is stable under sufficiently small bounded perturbations. Therefore $\sigma(A)$ is closed and bounded.

\textbf{Reality of the spectrum.}
Let $\lambda=a+ib$ with $b\neq 0$. Since $A$ is self-adjoint, $A-aI$ is self-adjoint. For every $x$,
\[
\|(A-\lambda I)x\|^2
=\|(A-aI)x-ibx\|^2
=\|(A-aI)x\|^2+b^2\|x\|^2.
\]
The cross terms cancel because $\langle (A-aI)x,x\rangle$ is real. Hence
\[
\|(A-\lambda I)x\|\geq |b|\|x\|.
\]
Thus $A-\lambda I$ is injective and has closed range.

To show the range is all of $H$, use
\[
\operatorname{Ran}(A-\lambda I)^\perp=\ker((A-\lambda I)^*)=\ker(A-\overline{\lambda}I).
\]
The same estimate with $\overline{\lambda}$ shows that this kernel is $\{0\}$. Therefore the range is dense; being both dense and closed, it is the whole space. The inverse exists and is bounded by the estimate above. Hence every nonreal $\lambda$ lies in the resolvent set, and
\[
\sigma(A)\subset\mathbb R.
\]
Combining the two parts, $\sigma(A)$ is a bounded closed subset of the real line.
\end{solution}

\subsection*{Team}

\begin{exercise}
\label{analysis:2016:7}
Let $D \subset \mathbb{R}^d, d \geq 2$ be a compact convex set with smooth boundary $\partial D$ so that the origin belongs to the interior of $D$. For every $x \in \partial D$ let $\alpha(x) \in (0, \infty)$ be the angle between the position vector $x$ and the outer normal vector $n(x)$. Let $\omega_d$ be the surface area of the unit sphere in $\mathbb{R}^d$. Compute:
\[
\frac{1}{\omega_d} \int_{\partial D} \frac{\cos(\alpha(x))}{|x|^{d-1}} d\sigma(x)
\]
where $d\sigma$ denotes the surface measure on $\partial D$.
\end{exercise}

\begin{solution}
The answer is $1$.

Let
\[
X(x)=\frac{x}{|x|^d},\qquad x\neq 0.
\]
Then
\[
X(x)\cdot n(x)=\frac{x\cdot n(x)}{|x|^d}
=\frac{\cos(\alpha(x))}{|x|^{d-1}},
\]
because $\cos(\alpha(x))=(x\cdot n(x))/|x|$.

The vector field $X$ is divergence-free away from the origin:
\[
\operatorname{div}\left(\frac{x}{|x|^d}\right)=0,\qquad x\neq 0.
\]
Since $0$ lies in the interior of $D$, remove a small ball $B_\rho(0)\subset D$ and apply the divergence theorem to $D\setminus B_\rho(0)$:
\[
0=\int_{\partial D} X\cdot n\,d\sigma
\;+\;\int_{\partial B_\rho(0)}X\cdot n_{D\setminus B_\rho}\,d\sigma .
\]
On the inner boundary, the outward normal for $D\setminus B_\rho(0)$ is $-\frac{x}{|x|}$, so
\[
X\cdot n_{D\setminus B_\rho}
=-\frac{1}{\rho^{d-1}}.
\]
Therefore
\[
\int_{\partial B_\rho(0)}X\cdot n_{D\setminus B_\rho}\,d\sigma
=-\omega_d.
\]
Hence
\[
\int_{\partial D}\frac{\cos(\alpha(x))}{|x|^{d-1}}\,d\sigma(x)
=\int_{\partial D}X\cdot n\,d\sigma
=\omega_d.
\]
Dividing by $\omega_d$ gives the required value:
\[
\frac{1}{\omega_d}\int_{\partial D}
\frac{\cos(\alpha(x))}{|x|^{d-1}}\,d\sigma(x)=1.
\]
\end{solution}

\begin{exercise}
\label{analysis:2016:8}
Let $p > 0$ and suppose $f_n, f \in L^p[0,1]$ and $\|f_n - f\|_p = (\int_0^1 |f_n(x) - f(x)|^p dx)^{1/p} \to 0$ as $n \to \infty$.
\begin{itemize}
    \item[a)] Show that for every $\epsilon > 0$
    \[
    \lim_{n \to \infty} m(\{x \in [0, 1] \mid |f_n(x) - f(x)| > \epsilon\}) = 0.
    \]
    Here $m$ is the Lebesgue measure.
    \item[b)] Show that there exists a subsequence $f_{n_j}$ such that $f_{n_j}(x) \to f(x)$ for almost every $x \in [0, 1]$.
\end{itemize}
\end{exercise}

\begin{solution}
Let
\[
g_n(x)=|f_n(x)-f(x)|^p.
\]
The hypothesis says
\[
\int_0^1 g_n(x)\,dx\to 0.
\]
This argument works for every $p>0$; no convexity of the $L^p$ norm is needed.

\textbf{Part (a).}
For every $\epsilon>0$,
\[
\{|f_n-f|>\epsilon\}=\{g_n>\epsilon^p\}.
\]
By Chebyshev's inequality,
\[
m(\{|f_n-f|>\epsilon\})
\leq \epsilon^{-p}\int_0^1 |f_n-f|^p\,dx.
\]
The right-hand side tends to $0$, proving convergence in measure.

\textbf{Part (b).}
Choose a subsequence $f_{n_j}$ so fast that
\[
\int_0^1 |f_{n_j}-f|^p\,dx<2^{-2j}.
\]
Then, by part (a) with $\epsilon_j=2^{-j/p}$,
\[
m\left(\{|f_{n_j}-f|>2^{-j/p}\}\right)
\leq 2^j\int_0^1 |f_{n_j}-f|^p\,dx
<2^{-j}.
\]
The series of these measures is finite. By the Borel--Cantelli lemma, for almost every $x$ only finitely many of the events
\[
\{|f_{n_j}(x)-f(x)|>2^{-j/p}\}
\]
occur. Hence for almost every $x$ there exists $j_0(x)$ such that for all $j\geq j_0(x)$,
\[
|f_{n_j}(x)-f(x)|\leq 2^{-j/p}\to 0.
\]
Thus $f_{n_j}(x)\to f(x)$ almost everywhere.
\end{solution}

\begin{exercise}
\label{analysis:2016:9}
\begin{itemize}
    \item[1)] Let $f$ be a holomorphic function on the unit disk $D = \{ z \in \mathbb{C} \mid |z| < 1 \}$ except 0. Assume $f \in L^2(D)$, i.e., $\int_D |f(z)|^2 dz d\bar{z} < \infty$, then 0 is a removable singularity.
    \item[2)] Let $f_n$ be a sequence of holomorphic functions over a domain $\Omega \subset \mathbb{C}$ converging to $f$ uniformly on any compact subset of $\Omega$, does the sequence of its derivatives $f_n'$ also have this property?
\end{itemize}
\end{exercise}

\begin{solution}
\textbf{Part 1.}
Write the Laurent expansion of $f$ in the punctured disk:
\[
f(z)=\sum_{n=-\infty}^{\infty}a_n z^n,\qquad 0<|z|<1.
\]
Fix $0<r<R<1$. Orthogonality of the exponentials on circles gives
\[
\int_{r<|z|<R}|f(z)|^2\,dA(z)
=2\pi\sum_{n=-\infty}^{\infty}|a_n|^2
\int_r^R \rho^{2n+1}\,d\rho.
\]
Here the identity is first applied on the fixed annulus $r<|z|<R$, where the Laurent series converges uniformly on compact subannuli, and then follows by Parseval on each circle. If some negative coefficient $a_n$ with $n<0$ were nonzero, then the corresponding radial integral would diverge as $r\to 0$: for $n=-1$ it is logarithmic, and for $n\leq -2$ it diverges as a negative power of $r$. This contradicts $f\in L^2(D)$.

Therefore all negative Laurent coefficients vanish. The Laurent series is actually a Taylor series,
\[
f(z)=\sum_{n=0}^{\infty}a_n z^n,
\]
so $f$ extends holomorphically to $0$. Thus the singularity is removable.

\textbf{Part 2.}
Yes. Let $K\Subset\Omega$ be compact. Choose a slightly larger compact set $K'\Subset\Omega$ and a radius $r>0$ such that every circle $\{|\zeta-z|=r\}$ with $z\in K$ is contained in $K'$. By Cauchy's integral formula,
\[
f_n'(z)-f'(z)
=\frac{1}{2\pi i}\int_{|\zeta-z|=r}
\frac{f_n(\zeta)-f(\zeta)}{(\zeta-z)^2}\,d\zeta.
\]
Hence
\[
|f_n'(z)-f'(z)|
\leq \frac{1}{r}\sup_{\zeta\in K'}|f_n(\zeta)-f(\zeta)|.
\]
The supremum on $K'$ tends to $0$ by compact-uniform convergence. Therefore $f_n'\to f'$ uniformly on $K$. Since $K$ was arbitrary, the derivatives converge uniformly on compact subsets of $\Omega$.
\end{solution}

\begin{exercise}
\label{analysis:2016:10}
Consider the torus $T^2 = \mathbb{C} / \Lambda, \Lambda = \{ m + in \mid m, n \in \mathbb{Z} \}$, i.e., $z_1, z_2 \in \mathbb{C}$ are equivalent if and only if there are integers $m, n$ such that $z_2 = z_1 + m + in$ and $T^2$ is the space of equivalence classes. Show that the group of holomorphic automorphisms of $T^2$ is $SL(2, \mathbb{Z})$ of $2 \times 2$ integer matrices of determinant 1.
\end{exercise}

\begin{solution}
As written, the statement is not correct for the complex torus
\[
T^2=\mathbb C/(\mathbb Z+i\mathbb Z).
\]
The group $SL(2,\mathbb Z)$ is the orientation-preserving mapping class group of the underlying smooth torus, not the group of holomorphic automorphisms for this fixed complex structure.

There are two quick obstructions.

First, every translation
\[
\tau_b([z])=[z+b],\qquad b\in T^2,
\]
is a holomorphic automorphism. Thus the full holomorphic automorphism group is uncountable, while $SL(2,\mathbb Z)$ is countable.

Second, if we restrict to automorphisms fixing the origin, a holomorphic lift has the form
\[
F(z)=az
\]
with $a\Lambda=\Lambda$. For the square lattice $\Lambda=\mathbb Z+i\mathbb Z$, this forces
\[
a\in\{\pm 1,\pm i\}.
\]
Equivalently, viewing a real torus map as an integer matrix
\[
A=\begin{pmatrix}p&q\\ r&s\end{pmatrix},
\]
holomorphicity means that $A$ commutes with the complex structure
\[
J=\begin{pmatrix}0&-1\\ 1&0\end{pmatrix}.
\]
This gives $s=p$ and $r=-q$, so
\[
A=\begin{pmatrix}p&q\\ -q&p\end{pmatrix}.
\]
The determinant condition is
\[
p^2+q^2=1,
\]
and the integer solutions are exactly
\[
(p,q)=(\pm 1,0),(0,\pm 1),
\]
corresponding to multiplication by $\pm1$ and $\pm i$.

Thus the correct holomorphic automorphism group is
\[
\operatorname{Aut}_{hol}(T^2)\cong T^2\rtimes\{\pm1,\pm i\}.
\]
If the intended problem was about orientation-preserving diffeomorphism classes of the real torus, then $SL(2,\mathbb Z)$ is the correct answer; but it is not the holomorphic automorphism group of the fixed complex torus stated here.
\end{solution}

\begin{exercise}
\label{analysis:2016:11}
Let $\{e_n\}$ be an orthonormal basis of $l_2$ of square integrable functions over a circle. Let $A: l_2 \to l_2$ be a linear operator defined by $A e_1 = 0, A e_n = \frac{e_{n-1}}{n-1}$ for $n > 1$. Show that $A$ is a compact operator and $A$ has no eigenvectors. What is the spectrum of $A$?
\end{exercise}

\begin{solution}
The compactness and spectrum can be computed explicitly. However, the sentence ``$A$ has no eigenvectors'' is false as written, because
\[
Ae_1=0.
\]
Thus $e_1$ is an eigenvector with eigenvalue $0$. The correct statement is that $A$ has no nonzero eigenvalues.

\textbf{Compactness.}
The operator is a weighted backward shift:
\[
Ae_n=w_n e_{n-1},\qquad n>1,\qquad w_n=\frac{1}{n-1},
\]
and $Ae_1=0$. Let $A_N$ be the finite-rank truncation defined by
\[
A_Ne_n=
\begin{cases}
Ae_n,& n\leq N,\\
0,& n>N.
\end{cases}
\]
Then
\[
\|A-A_N\|=\sup_{n>N} w_n=\frac{1}{N}.
\]
Thus $A$ is the operator-norm limit of finite-rank operators, so $A$ is compact.

\textbf{Eigenvalues.}
Let $x=\sum_{n\geq1}c_ne_n$ and suppose $Ax=\lambda x$. Comparing coefficients of $e_k$ gives
\[
\frac{c_{k+1}}{k}=\lambda c_k,\qquad k\geq1.
\]
Hence
\[
c_{k+1}=\lambda k c_k
\]
and therefore
\[
c_n=\lambda^{n-1}(n-1)!c_1.
\]
If $\lambda\neq0$ and $c_1\neq0$, the coefficients do not form a square-summable sequence. If $c_1=0$, all $c_n=0$. Hence there are no nonzero eigenvalues. For $\lambda=0$, the equation gives $c_n=0$ for $n\geq2$, while $c_1$ is arbitrary, so
\[
\ker A=\operatorname{span}\{e_1\}.
\]

\textbf{Spectrum.}
For $m>n$,
\[
A^n e_m=
\frac{1}{(m-1)(m-2)\cdots(m-n)}e_{m-n}.
\]
The largest coefficient occurs at $m=n+1$, so
\[
\|A^n\|=\frac{1}{n!}.
\]
The spectral radius formula gives
\[
r(A)=\lim_{n\to\infty}\|A^n\|^{1/n}
=\lim_{n\to\infty}(n!)^{-1/n}=0.
\]
Since the spectrum is nonempty and compact, and every spectral value has modulus at most $r(A)$, we get
\[
\sigma(A)=\{0\}.
\]
Thus the accurate conclusion is: $A$ is compact, its spectrum is $\{0\}$, $0$ is an eigenvalue with eigenvector $e_1$, and there are no nonzero eigenvalues.
\end{solution}

\begin{exercise}
\label{analysis:2016:12}
If $M = [0, 1]$ is the unit interval, the heat kernel on $M$ can be written
\[
p(x, y, t) = \sum_k \phi_k(x) \phi_k(y) e^{\lambda_k t},
\]
where $\{\lambda_k\}$ is an enumeration of the eigenvalues of $\Delta = \frac{d^2}{dx^2}$ on $M$ and $\{\phi_k\}$ are the corresponding eigenfunctions which vanish on $\partial M$.
\begin{itemize}
    \item[i)] Calculate $\{\lambda_k\}$ and the corresponding eigenfunctions.
    \item[ii)] Prove that $|p(x, y, t)| \leq C t^{-1/2}$, for all $x, y$, and $0 < t < 1$.
    \item[iii)] What is the exponential rate of decay of $p(x, y, t)$ as $t \to \infty$, i.e., compute:
    \[
    \lim_{t \to \infty} \log(p(x, y, t)).
    \]
\end{itemize}
\end{exercise}

\begin{solution}
We impose the Dirichlet boundary condition $\phi(0)=\phi(1)=0$. The eigenvalue problem is
\[
\phi''(x)=\lambda\phi(x),\qquad \phi(0)=\phi(1)=0.
\]
The nonzero solutions are
\[
\lambda_k=-k^2\pi^2,\qquad
\phi_k(x)=\sqrt2\sin(k\pi x),\qquad k=1,2,\dots,
\]
where the factor $\sqrt2$ normalizes the eigenfunctions in $L^2[0,1]$.

Thus
\[
p(x,y,t)=2\sum_{k=1}^\infty
\sin(k\pi x)\sin(k\pi y)e^{-k^2\pi^2t}.
\]

\textbf{The small-time bound.}
Since $|\sin(k\pi x)\sin(k\pi y)|\leq1$,
\[
|p(x,y,t)|\leq 2\sum_{k=1}^\infty e^{-k^2\pi^2t}.
\]
For $0<t<1$, compare the sum with a Gaussian integral:
\[
\sum_{k=1}^\infty e^{-k^2\pi^2t}
\leq \int_0^\infty e^{-\pi^2t s^2}\,ds
=\frac{1}{2\sqrt{\pi t}}.
\]
This gives
\[
|p(x,y,t)|\leq C t^{-1/2}
\]
for a universal constant $C$.

\textbf{Large-time decay.}
For $x,y\in(0,1)$, the first eigenfunction term dominates:
\[
p(x,y,t)
=2\sin(\pi x)\sin(\pi y)e^{-\pi^2t}
+O(e^{-4\pi^2t}).
\]
Therefore the exponential decay rate is
\[
\lim_{t\to\infty}\frac{1}{t}\log p(x,y,t)=-\pi^2.
\]
If one reads the displayed expression in the problem literally, without the factor $1/t$, then
\[
\lim_{t\to\infty}\log p(x,y,t)=-\infty
\]
for interior $x,y$. At boundary points the Dirichlet heat kernel is identically zero, so the logarithm is not finite there.
\end{solution}
